\chapter[20]{\texorpdfstring{Compactness of the Space of $\kappa$-Solutions}{Compactness of the Space of Kappa-Solutions}}
\label{ch:chow-iii-compactness-kappa-solutions}
\dcref{ch20}
\source{chow-et-al-ricci-flow-part-iii:page-0144,page-0145}

\section{Asymptotic scalar curvature and volume ratios}

\begin{theorem}[$\kappa$-solutions with Harnack have ASCR $=\infty$ and AVR $=0$]
\label{thm:chow-iii-kappa-harnack-ascr-avr}
\source{chow-et-al-ricci-flow-part-iii:page-0145,page-0146,page-0147,page-0148,page-0149,page-0150}
\uses{def:chow-iii-kappa-harnack, def:chow-iii-asymptotic-ratios, thm:chow-iii-point-picking-infinite-ascr, thm:chow-iii-dimension-reduction-one, thm:chow-iii-2d-kappa}
\dcref{ch20:20.1}
\group{Asymptotic Geometry}


For every complete noncompact $n$-dimensional $\kappa$-solution with Harnack,
$n\ge2$, one has
\[
 \operatorname{ASCR}(g(t))=\infty,
 \qquad \operatorname{AVR}(g(t))=0
\]
for every $t\le0$.
\end{theorem}
\begin{proof}
\sketch Argue by induction on $n$.  There are no noncompact two-dimensional
$\kappa$-solutions with Harnack.  At a fixed time $t_0$, failure of the
conclusion leaves three cases.

If $\operatorname{ASCR}=\infty$ and $\operatorname{AVR}>0$, spatial point
picking and curvature rescaling give a nonflat pointed limit containing a
line.  It splits as $\mathbb R\times W^{n-1}$, where $W$ is a
$\kappa$-solution.  Bishop--Gromov comparison passes positive AVR to the
limit and then to $W$, contradicting the induction hypothesis.

If $0<\operatorname{ASCR}<\infty$, a blow-down on annuli converges both
smoothly to a nonflat Ricci flow and in the Gromov--Hausdorff sense to the
regular part of the asymptotic cone
\[
 dr^2+r^2g_W.
\]
All radial two-planes lie in the null space of its curvature operator.
Hamilton's strong maximum principle makes that null space parallel; its
parallel translates span $\Lambda^2$, so the limit is flat, a contradiction.

Finally, if $\operatorname{ASCR}=0$, the Petrunin--Tuschmann gap theorem and
the strong maximum principle split the universal cover as
$\mathbb R^{n-2}\times\Sigma^2$, where $\Sigma$ is a two-dimensional
$\kappa$-solution.  Such a solution is round, which is incompatible with
zero ASCR.  Thus none of the three cases occurs.
\end{proof}

\begin{corollary}[Any ancient solution with bounded $\Rm\ge0$ must have AVR $=0$]
\label{cor:chow-iii-ancient-bounded-nonnegative-avr-zero}
\source{chow-et-al-ricci-flow-part-iii:page-0145,page-0146}
\uses{thm:chow-iii-kappa-harnack-ascr-avr}
\dcref{ch20:20.2}
\group{Asymptotic Geometry}


Let $(\mathcal M^n,g(t))$, $t\le0$, be complete, noncompact, nonflat, and
ancient, with nonnegative curvature operator and
\[
 \sup_{\mathcal M\times(-\infty,0]}|\Rm|<\infty.
\]
Then $\operatorname{AVR}(g(t))=0$ for every $t\le0$.
\end{corollary}
\begin{proof}
\sketch If $\omega_n\operatorname{AVR}(g(t_0))=\kappa_0>0$, Bishop--Gromov
comparison makes $g(t_0)$ $\kappa_0$-noncollapsed at every scale.  The proof
of Theorem~\ref{thm:chow-iii-kappa-harnack-ascr-avr} uses only this
time-slice noncollapsing to force zero AVR, a contradiction.
\end{proof}

\begin{remark}[Kahler analogue]
\label{rem:chow-iii-kahler-ancient-avr-zero}
\source{chow-et-al-ricci-flow-part-iii:page-0146}
\dcref{ch20:20.3}
\group{Asymptotic Geometry}

The analogous statement holds for a nonflat ancient Kahler--Ricci flow with
bounded nonnegative bisectional curvature: its asymptotic volume ratio is
zero at every time.
\end{remark}

\section{\texorpdfstring{Almost $\kappa$-solutions}{Almost Kappa-solutions}}

\begin{proposition}[Almost ancient solutions with bounded $\Rm\ge0$ are collapsed at large scales]
\label{prop:chow-iii-almost-ancient-large-scale-collapse}
\source{chow-et-al-ricci-flow-part-iii:page-0151,page-0152}
\uses{cor:chow-iii-ancient-bounded-nonnegative-avr-zero, prop:chow-iii-cgt-injectivity}
\dcref{ch20:20.4}
\group{Almost Kappa-Solutions}


For every $\varepsilon>0$ there is $A<\infty$ with the following property.
Suppose $(\mathcal M_k^n,g_k(t))$, $t\in[t_k,0]$, have nonnegative curvature
operator, and there are $x_k\in\mathcal M_k$, $r_k>0$ such that
$B_{g_k(0)}(x_k,r_k)\Subset\mathcal M_k$,
\[
 \frac12R_{g_k}(x,t)\le R_{g_k}(x_k,0)=Q_k
\]
on $B_{g_k(0)}(x_k,r_k)\times(t_k,0]$, and
\[
 t_kQ_k\longrightarrow-\infty,
 \qquad r_k^2Q_k\longrightarrow\infty.
\]
Then, for all sufficiently large $k$,
\[
 \frac{\Vol_{g_k(0)}B_{g_k(0)}(x_k,AQ_k^{-1/2})}
 {(AQ_k^{-1/2})^n}\le\varepsilon.
\]
\end{proposition}
\begin{proof}
\sketch Otherwise choose $A_k\to\infty$ for which the displayed volume ratios have a
fixed positive lower bound.  After rescaling by $Q_k$, Bishop--Gromov and the
curvature bound give a uniform injectivity radius at $x_k$.  Local Hamilton
compactness yields a complete, nonflat, bounded-curvature ancient limit with
$\Rm\ge0$.  The volume lower bounds pass to the limit and give positive AVR,
contradicting Corollary~\ref{cor:chow-iii-ancient-bounded-nonnegative-avr-zero}.
\end{proof}

\begin{remark}[Interpretation of almost ancient]
\label{rem:chow-iii-almost-ancient-rescaling}
\source{chow-et-al-ricci-flow-part-iii:page-0151}
\uses{prop:chow-iii-almost-ancient-large-scale-collapse}
\dcref{ch20:20.5}
\group{Almost Kappa-Solutions}


For $\widetilde g_k(t)=Q_kg_k(Q_k^{-1}t)$, the hypotheses say that the
backward time intervals tend to $(-\infty,0]$ and that $R\le2$ on balls whose
radii tend to infinity.  Thus the sequence becomes ancient and globally
bounded in the pointed limit.
\end{remark}

\begin{proposition}[Curvature estimate in noncollapsed balls]
\label{prop:chow-iii-curvature-estimate-noncollapsed-balls}
\source{chow-et-al-ricci-flow-part-iii:page-0152,page-0153,page-0154,page-0155,page-0156,page-0157}
\uses{prop:chow-iii-almost-ancient-large-scale-collapse, prop:chow-iii-parabolic-cylinder-curvature-control, thm:chow-iii-changing-distances}
\dcref{ch20:20.6}
\group{Almost Kappa-Solutions}


For every $\omega>0$ there are $B(\omega),C(\omega)<\infty$ and
$\tau_0(\omega)>0$ as follows.  Let $(\mathcal M^n,g(t))$,
$t\in[t_0,0]$, be a not necessarily complete Ricci flow with $\Rm\ge0$, and
assume $B_{g(0)}(x_0,r_0)\Subset\mathcal M$.

If
\[
 \Vol_{g(t)}B_{g(t)}(x_0,r_0)\ge\omega r_0^n
 \quad(t_0\le t\le0),
\]
then, whenever $d_{g(t)}(x,x_0)\le r_0/4$,
\[
 R(x,t)\le C r_0^{-2}+B(t-t_0)^{-1}.
\]
If instead only
\[
 \Vol_{g(0)}B_{g(0)}(x_0,r_0)\ge\omega r_0^n
\]
is assumed, then for $-t_0\ge\tau_0r_0^2$ and
$d_{g(t)}(x,x_0)\le r_0/4$,
\[
 R(x,t)\le C r_0^{-2}+B(t+\tau_0r_0^2)^{-1},
 \qquad -\tau_0r_0^2<t\le0.
\]
\end{proposition}
\begin{proof}
\sketch Scale to $r_0=1$.  If the first assertion fails, parabolic point picking
produces points $(y_k,t_k^*)$ with curvature scales tending to infinity and
uniform curvature bounds on increasingly large backward parabolic
cylinders.  Bishop--Gromov comparison transfers the assumed unit-ball volume
lower bound to these curvature-scale balls.  Proposition
\ref{prop:chow-iii-almost-ancient-large-scale-collapse} forces their volume
ratios to tend to zero, a contradiction.

For the second assertion, follow the fixed final-time ball backward until its
points first reach distance $19/80$ from $x_0$.  The first assertion supplies
a curvature bound during this interval.  The changing-distance estimate then
shows, after choosing $\tau_0(\omega)$ sufficiently small, that this first
time occurs no later than $-\tau_0$.  Ricci nonnegativity and Bishop--Gromov
comparison propagate a definite volume lower bound throughout
$[-\tau_0,0]$, so the first assertion applies there.
\end{proof}

\begin{remark}[Independence from the short-time parameter]
\label{rem:chow-iii-curvature-estimate-time-independent}
\source{chow-et-al-ricci-flow-part-iii:page-0153}
\dcref{ch20:20.7}
\group{Almost Kappa-Solutions}

The first, time-interval-noncollapsed assertion of Proposition
\ref{prop:chow-iii-curvature-estimate-noncollapsed-balls} is independent of
the choice of $\tau_0$.
\end{remark}

\section{Compactness modulo scaling}

\begin{definition}[Precompact and compact modulo scaling]
\label{def:chow-iii-compact-modulo-scaling}
\source{chow-et-al-ricci-flow-part-iii:page-0157,page-0158}
\dcref{ch20:20.21--20.22}
\group{Compactness}

For a sequence of ancient flows with positive scalar curvature and chosen
space-time points $(x_k,t_k)$, define the curvature-normalized flows
\[
 \bar g_k(t)=R_{g_k}(x_k,t_k)
 g_k\!\left(t_k+\frac{t}{R_{g_k}(x_k,t_k)}\right),
 \qquad R_{\bar g_k}(x_k,0)=1.
\]
A collection is \emph{precompact modulo scaling} if every such normalized
sequence has a pointed smooth convergent subsequence.  It is \emph{compact
modulo scaling} if every limit remains in the collection.
\end{definition}

\begin{remark}[Positive scalar curvature or Ricci flatness]
\label{rem:chow-iii-ancient-positive-or-ricci-flat}
\source{chow-et-al-ricci-flow-part-iii:page-0158}
\dcref{ch20:20.8}
\group{Compactness}

An ancient solution with bounded curvature either has positive scalar
curvature or is Ricci flat.
\end{remark}

\begin{theorem}[Compactness of $\kappa$-solutions with Harnack]
\label{thm:chow-iii-kappa-harnack-compactness}
\source{chow-et-al-ricci-flow-part-iii:page-0158,page-0160,page-0161,page-0162,page-0163,page-0164,page-0165,page-0166,page-0167,page-0168,page-0169,page-0170}
\uses{def:chow-iii-compact-modulo-scaling, thm:chow-iii-kappa-harnack-ascr-avr, prop:chow-iii-almost-ancient-large-scale-collapse, prop:chow-iii-curvature-estimate-noncollapsed-balls, lem:chow-iii-bounded-curvature-bounded-distance, clm:chow-iii-unit-ball-volume-lower-bound}
\dcref{ch20:20.9}
\group{Compactness}


For every $\kappa>0$ and $n\ge2$, the collection of $n$-dimensional
$\kappa$-solutions with Harnack is compact modulo scaling.
\end{theorem}
\begin{proof}
\sketch Normalize an arbitrary pointed sequence so that $R(x_k,0)=1$.  One proof
sets
\[
 A_k=\sup_z R(z,0)d(x_k,z)^2.
\]
If $A_k\le1$, Theorem~\ref{thm:chow-iii-kappa-harnack-ascr-avr} makes each
manifold compact.  A second rescaling at a hypothetical maximum-curvature
point would produce, at bounded distance, a smooth limit with scalar
curvature both zero and positive, contrary to the strong maximum principle.
Thus curvature is uniformly bounded.

If $A_k>1$, choose the closest $z_k$ with
$R(z_k,0)d(x_k,z_k)^2=1$.  Point picking, noncollapsing, and Proposition
\ref{prop:chow-iii-almost-ancient-large-scale-collapse} first bound curvature
relative to $R(z_k,0)$ on a ball about $z_k$.  The integrated Harnack estimate
and Shi estimates then bound $R(z_k,0)$ absolutely.  Lemma
\ref{lem:chow-iii-bounded-curvature-bounded-distance} extends this to every
fixed-radius ball about $x_k$.

Alternatively, Claim~\ref{clm:chow-iii-unit-ball-volume-lower-bound} gives a
uniform volume lower bound at $x_k$, and Proposition
\ref{prop:chow-iii-curvature-estimate-noncollapsed-balls} directly bounds
curvature on every fixed-radius ball.  Either route gives the local curvature
and injectivity bounds needed for Hamilton compactness.  Completeness,
noncollapsing, nonnegative curvature operator, trace Harnack, and the
normalization pass to the smooth limit.
\end{proof}

\begin{corollary}[Compactness of three-dimensional $\kappa$-solutions]
\label{cor:chow-iii-3d-kappa-compactness}
\source{chow-et-al-ricci-flow-part-iii:page-0158}
\uses{thm:chow-iii-kappa-harnack-compactness, prop:chow-iii-3d-harnack-equivalence}
\dcref{ch20:20.10}
\group{Compactness}


For fixed $\kappa>0$, the collection of three-dimensional $\kappa$-solutions
is compact modulo scaling.
\end{corollary}

\begin{corollary}[Improved compactness of three-dimensional $\kappa$-solutions]
\label{cor:chow-iii-3d-kappa-improved-compactness}
\source{chow-et-al-ricci-flow-part-iii:page-0158,page-0159}
\uses{cor:chow-iii-3d-kappa-compactness, thm:chow-iii-kappa-gap}
\dcref{ch20:20.11}
\group{Compactness}


The union, over all $\kappa>0$, of the noncompact three-dimensional
$\kappa$-solutions is compact modulo scaling.  The corresponding union of
nonspherical-space-form three-dimensional $\kappa$-solutions is precompact
modulo scaling, although a limit may be spherical.
\end{corollary}
\begin{proof}
\sketch The three-dimensional gap theorem supplies one uniform $\kappa_0>0$ for all
members of either union.  Apply Corollary
\ref{cor:chow-iii-3d-kappa-compactness}.  Noncompactness is preserved in the
first family; in the second family only precompactness is asserted because a
spherical limit is possible.
\end{proof}

\begin{example}[Pointed limits of model $\kappa$-solutions]
\label{ex:chow-iii-model-kappa-pointed-limits}
\source{chow-et-al-ricci-flow-part-iii:page-0159}
\uses{ex:chow-iii-bryant-soliton, thm:chow-iii-perelman-sphere}
\dcref{ch20:20.12}
\group{Compactness}


Pointed curvature-normalized Bryant solitons can limit only to the Bryant
soliton or the round cylinder.  The expected limits of Perelman's spherical
$\kappa$-solution are the shrinking round sphere, the solution itself, the
Bryant soliton, and the round cylinder, according to the chosen times and the
location of the basepoints relative to the tips.
\end{example}

\begin{remark}[Exclusion of conical degeneration]
\label{rem:chow-iii-kappa-compactness-no-conical-degeneration}
\source{chow-et-al-ricci-flow-part-iii:page-0160}
\uses{thm:chow-iii-kappa-harnack-compactness}
\dcref{ch20:20.13}
\group{Compactness}


The compactness theorem rules out curvature blow-up a uniformly bounded
distance from basepoints of bounded curvature; in particular, a sequence
cannot develop a singular Euclidean-cone pointed limit at such a point.
\end{remark}

\begin{remark}[Rescaling at a high-curvature point]
\label{rem:chow-iii-rescaling-bounded-distance}
\source{chow-et-al-ricci-flow-part-iii:page-0162}
\dcref{ch20:20.14}
\group{Compactness}

If a closed manifold has $R(x)=1$ and $R(y)d(x,y)^2\le C$, then for
$\widetilde g=R(w)g$ one has
\[
 R_{\widetilde g}(x)=R(w)^{-1},\qquad
 R_{\widetilde g}(w)=1,\qquad
 d_{\widetilde g}(x,w)\le\sqrt C.
\]
Thus unbounded $R(w)$ would put a zero-curvature point and a unit-curvature
point at bounded distance in a pointed limit.
\end{remark}

\begin{lemma}[Bounded curvature at bounded distance]
\label{lem:chow-iii-bounded-curvature-bounded-distance}
\source{chow-et-al-ricci-flow-part-iii:page-0167,page-0168}
\uses{prop:chow-iii-almost-ancient-large-scale-collapse, lem:chow-iii-point-picking-unbounded-change}
\dcref{ch20:20.15}
\group{Compactness}


Let $(\mathcal M_k,g_k(t),x_k)$, $t\in(-\alpha_k,0]$, be complete Ricci
flows with $\Rm\ge0$, $\partial_tR\ge0$, and suppose
\[
 R(x_k,0)=1,
 \qquad \Vol_{g_k(0)}B_{g_k(0)}(x_k,r_0)\ge v_0>0.
\]
For every $r>0$, either
\[
 \alpha_k\sup_{B_{g_k(0)}(x_k,r)}R(\cdot,0)\le C_0
\]
uniformly in $k$, or
\[
 \sup_{B_{g_k(0)}(x_k,r)}R(\cdot,0)\le C_1
\]
uniformly in $k$.
\end{lemma}
\begin{proof}
\sketch If both alternatives fail, point picking supplies $w_k$ and $s_k$ with
$Q_ks_k^2\to\infty$, $\alpha_kQ_k\to\infty$, and $R\le2Q_k$ on
$B(w_k,s_k)\times(-\alpha_k,0]$.  Proposition
\ref{prop:chow-iii-almost-ancient-large-scale-collapse} makes the volume
ratio at a fixed multiple of $Q_k^{-1/2}$ arbitrarily small.  The ball about
$x_k$ with volume at least $v_0$ lies in a fixed-radius ball about $w_k$;
Bishop--Gromov comparison gives the opposite positive lower bound.
\end{proof}

\begin{claim}[Volume noncollapsing from curvature bounded at the center]
\label{clm:chow-iii-unit-ball-volume-lower-bound}
\source{chow-et-al-ricci-flow-part-iii:page-0169,page-0170}
\uses{prop:chow-iii-curvature-estimate-noncollapsed-balls}
\dcref{ch20:20.16}
\group{Compactness}


For curvature-normalized $n$-dimensional $\kappa$-solutions with Harnack,
there is $c_0(\kappa,n)>0$ such that
\[
 \Vol_{g_k(0)}B_{g_k(0)}(x_k,1)\ge c_0.
\]
\end{claim}
\begin{proof}
\sketch If the volumes tend to zero, choose $\delta_k\downarrow0$ so that the volume
ratio at radius $\delta_k$ is $\omega_n/2$ and rescale by $\delta_k^{-2}$.
Proposition~\ref{prop:chow-iii-curvature-estimate-noncollapsed-balls} gives
curvature bounds on every fixed ball, hence a complete pointed limit.  Its
basepoint scalar curvature is zero, so the strong maximum principle makes it
flat.  Noncollapsing makes a complete flat limit Euclidean, contradicting
the unit-ball volume $\omega_n/2$.
\end{proof}

\section{Derivative estimates and applications}

\begin{theorem}[Derivative estimates for $\kappa$-solutions satisfying trace Harnack]
\label{thm:chow-iii-kappa-scaled-derivative-estimates}
\source{chow-et-al-ricci-flow-part-iii:page-0170,page-0171}
\uses{thm:chow-iii-kappa-harnack-compactness}
\dcref{ch20:20.17}
\group{Derivative Estimates}


For $\kappa>0$ and $n\ge2$, there is $\eta(\kappa,n)<\infty$ such that every
$n$-dimensional $\kappa$-solution with Harnack satisfies
\[
 |\nabla R|\le\eta R^{3/2},
 \qquad |\partial_tR|\le\eta R^2.
\]
\end{theorem}
\begin{proof}
\sketch Otherwise rescale a violating sequence at the chosen space-time points so
that $R(x_k,0)=1$.  One of $|\nabla R|(x_k,0)$ and
$|\partial_tR|(x_k,0)$ tends to infinity.  Theorem
\ref{thm:chow-iii-kappa-harnack-compactness} gives pointed smooth convergence,
under which both quantities converge to finite values, a contradiction.
\end{proof}

\begin{remark}[Estimates for derivatives of any order]
\label{rem:chow-iii-kappa-all-derivative-estimates}
\source{chow-et-al-ricci-flow-part-iii:page-0171}
\uses{thm:chow-iii-kappa-scaled-derivative-estimates}
\dcref{ch20:20.18}
\group{Derivative Estimates}


For every $\ell,m\ge0$ there is $\eta(\kappa,\ell,m,n)$ such that
\[
 |\partial_t^\ell\nabla^mR|
 \le\eta(\kappa,\ell,m,n)R^{1+\ell+m/2}.
\]
The same scale-invariant estimate holds for any ordering of space and time
derivatives of $\Rm$, $\Ric$, or $R$.
\end{remark}

\begin{corollary}[Universal scaled derivative bounds in dimension three]
\label{cor:chow-iii-3d-universal-scaled-derivatives}
\source{chow-et-al-ricci-flow-part-iii:page-0171,page-0172}
\uses{cor:chow-iii-3d-kappa-improved-compactness, thm:chow-iii-kappa-scaled-derivative-estimates}
\dcref{ch20:20.19}
\group{Derivative Estimates}


There is a universal $\eta_0<\infty$, independent of $\kappa$, such that
every three-dimensional $\kappa$-solution satisfies
\[
 |\nabla R|\le\eta_0R^{3/2},
 \qquad |\partial_tR|\le\eta_0R^2.
\]
More generally, constants $\eta(\ell,m)$ independent of $\kappa$ bound
$|\partial_t^\ell\nabla^m\Rm|$ by
$\eta(\ell,m)R^{1+\ell+m/2}$.
\end{corollary}

\begin{remark}[Derivative estimate fails on the cigar soliton]
\label{rem:chow-iii-cigar-scaled-derivative-failure}
\source{chow-et-al-ricci-flow-part-iii:page-0172}
\dcref{ch20:20.20}
\group{Derivative Estimates}

For the cigar metric $ds^2+\tanh^2s\,d\theta^2$ with
$R=4\operatorname{sech}^2s$, one has
$\sup R^{-3/2}|\nabla R|=\infty$.  Curvature rescaling about points tending
to infinity produces a pointed half-line rather than a smooth surface limit.
\end{remark}

\begin{corollary}[Quadratic scalar-curvature lower bound]
\label{cor:chow-iii-kappa-quadratic-curvature-lower-bound}
\source{chow-et-al-ricci-flow-part-iii:page-0172}
\uses{thm:chow-iii-kappa-scaled-derivative-estimates}
\dcref{ch20:20.21}
\group{Derivative Applications}


For a noncompact $\kappa$-solution with Harnack and any fixed $p$,
\[
 \liminf_{d_{g(t)}(x,p)\to\infty}R(x,t)d_{g(t)}(x,p)^2
 \ge\frac4{\eta^2}.
\]
\end{corollary}
\begin{proof}
\sketch Along a unit-speed minimizing geodesic from $p$ to $x$,
$\frac d{ds}R^{-1/2}\le\eta/2$.  Hence
\[
 R(x,t)\ge
 \left(R(p,t)^{-1/2}+\frac\eta2d_{g(t)}(x,p)\right)^{-2},
\]
and the assertion follows after multiplying by the squared distance.
\end{proof}

\begin{corollary}[Scalar curvature along steady-soliton flow lines]
\label{cor:chow-iii-steady-soliton-flowline-curvature-lower-bound}
\source{chow-et-al-ricci-flow-part-iii:page-0173}
\uses{thm:chow-iii-kappa-scaled-derivative-estimates}
\dcref{ch20:20.23}
\group{Derivative Applications}


Let a nonflat steady gradient Ricci soliton have $\Rm\ge0$ and be
$\kappa$-noncollapsed at all scales.  If
$\dot\sigma=-\nabla f$, then
\[
 R(\sigma(u))\ge
 \frac1{R(\sigma(0))^{-1}+\eta u}.
\]
In dimension three this implies $R(x)\ge c/(1+d(x,p))$.
\end{corollary}
\begin{proof}
\sketch The soliton satisfies trace Harnack, and Theorem
\ref{thm:chow-iii-kappa-scaled-derivative-estimates} gives
$\langle\nabla R,\nabla f\rangle\le\eta R^2$.  Therefore
$\frac d{du}R(\sigma(u))^{-1}\le\eta$, which integrates to the result.
\end{proof}

\begin{corollary}[Linear scalar-curvature upper bound for steady solitons]
\label{cor:chow-iii-3d-steady-soliton-curvature-upper-bound}
\source{chow-et-al-ricci-flow-part-iii:page-0173,page-0174}
\dcref{ch20:20.24}
\group{Derivative Applications}

For a three-dimensional nonflat steady gradient Ricci soliton with
$\Rm\ge0$ which is $\kappa$-noncollapsed at all scales, there is $C<\infty$
such that
\[
 R(x)\le\frac C{1+d(x,p)}.
\]
\end{corollary}

\begin{claim}[Optimistic conjecture in dimension two]
\label{conj:chow-iii-2d-ancient-classification}
\dcref{ch20:20.27}
\group{Classification Conjectures}
\source{chow-et-al-ricci-flow-part-iii:page-0174,page-0175}
Every two-dimensional ancient solution with bounded curvature is one of:
a quotient of $S^2$ or $\mathbb R^2$ with constant nonnegative curvature;
the cigar steady soliton; or the King--Rosenau ancient solution on $S^2$ or
$\mathbb{RP}^2$.
\end{claim}

\begin{claim}[Optimistic conjecture in dimension three]
\label{conj:chow-iii-3d-positive-kappa-classification}
\dcref{ch20:20.26}
\group{Classification Conjectures}
\source{chow-et-al-ricci-flow-part-iii:page-0174}
The only three-dimensional $\kappa$-solutions with positive sectional
curvature are round spherical space forms, the Bryant steady soliton, and
Perelman's spherical $\kappa$-solution on $S^3$ or $\mathbb{RP}^3$.
\end{claim}

\begin{example}[Problem: infinite asymptotic scalar curvature from derivatives]
\label{prob:chow-iii-kappa-derivatives-infinite-ascr}
\dcref{ch20:20.22}
\group{Derivative Applications}
\source{chow-et-al-ricci-flow-part-iii:page-0173}
Can the lower bound in Corollary
\ref{cor:chow-iii-kappa-quadratic-curvature-lower-bound} be strengthened to
\[
 \liminf_{d_{g(t)}(x,p)\to\infty}R(x,t)d_{g(t)}(x,p)^2=\infty?
\]
\end{example}

\begin{example}[Problem: rotationally symmetric positive $\kappa$-solutions]
\label{prob:chow-iii-rotational-kappa-bryant}
\dcref{ch20:20.25}
\group{Classification Problems}
\source{chow-et-al-ricci-flow-part-iii:page-0174}
Prove that every rotationally symmetric noncompact three-dimensional
$\kappa$-solution with positive sectional curvature is a constant multiple
of the Bryant soliton.
\end{example}
